% SPDX-License-Identifier: Apache-2.0
% covers: Spi
\datasheet{Spi}{An SPI master on AXI-Lite: one byte each way at a time, under a chip select a program holds.}

\dsfacts{\code{lib/parts/src/spi.rs}}{\code{txhdl\_parts::spi::Spi}}%
{\code{//docs:examples}}

\subsection*{Function}

\code{Spi} drives the four wires of an SPI link: \code{sclk}, the
clock; \code{mosi}, the byte going out; \code{miso}, the byte coming
in; and \code{cs\_n}, the chip select, low while a chip is held.

A program writes a byte to \code{data}, waits for the transfer to
finish, and reads \code{data} for the byte that arrived the other way.
SPI has no direction: the eight clocks that send a byte are the eight
that receive one, so every transfer is a swap, and a driver that wants
to read sends a byte it does not care about.

The chip select is a bit in \code{ctrl} that a program sets and
clears, not something a transfer does. Every SPI command is several
bytes under one unbroken select, so a master that dropped it between
bytes would start a new command each time.

\subsection*{Parameters}

None. The divider and the mode are registers rather than parameters,
so one instance can reach chips that want different clocks.

\subsection*{Register map}

Offsets from the peripheral's base. The part selects the word by bits
3 and 2 of the address and looks at no other bits.

\begin{center}\footnotesize
\begin{tabular}{@{}llp{0.5\textwidth}@{}}
\toprule
Offset & Word & Access \\
\midrule
\code{0x00} & \code{ctrl} & Read and write. Bits 7 to 0 the divider,
8 \code{cpol}, 9 \code{cpha}, 10 the chip select, 11 the interrupt
enable. \\
\code{0x04} & \code{data} & Written: the byte to send, which starts a
transfer. Read: the byte that came in. \\
\code{0x08} & \code{state} & Bit 0 a transfer is running, bit 1 one
has finished. Write a one to bit 1 to clear it. \\
\bottomrule
\end{tabular}
\end{center}

A half of a bit takes \code{div + 1} cycles, so the clock is the
system's over \code{2 * (div + 1)}.

\dstables{Spi}

\subsection*{Behaviour}

\begin{itemize}
\item A write to \code{data} while no transfer is running loads
\code{txb}, sets \code{busy}, and clears \code{tick}, \code{half} and
\code{count}. A write while one is running is answered and does
nothing, so a driver reads \code{state} first.
\item While \code{busy}, \code{tick} counts to \code{div}. The cycle
it reaches it, \code{half} turns over and \code{count} adds one.
\code{sclk} is \code{cpol} exclusive-or \code{half}, so it idles at
\code{cpol}.
\item The bit is taken on the leading turn when \code{cpha} is clear
and on the trailing turn when it is set. \code{rxb} shifts left and
takes \code{miso}.
\item \code{txb} shifts on the other turn, with one exception: when
\code{cpha} is set the first leading turn does not shift, because the
first bit is already on the wire and moving it would lose the last bit
of the byte. That is the turn at \code{count} zero.
\item After sixteen turns, eight bits, \code{busy} clears,
\code{half} returns to zero and \code{fired} is set. \code{irq} is
\code{fired} and the interrupt enable.
\item \code{mosi} is the top bit of \code{txb} and \code{cs\_n} is the
inverse of the select bit, both combinational.
\item A read is answered on \code{bus\_r} in the cycle it is taken, with
\code{OKAY}. A write is taken when its address phase, its beat and
room on \code{bus\_b} are all there.
\end{itemize}

\subsection*{Verification}

\begin{itemize}
\item \code{//lib/examples:ex\_spi} reads the chip's identity with
\code{0x9f} and four bytes from an address with the \code{0x0b} fast
read, against \code{FlashDevice}, and asserts both against what was
put in the model.
\item The build simulates \code{Spi}'s netlist against that run under
nvc and Verilator: \code{//docs:spi\_sim\_spi\_tb\_test} and
\code{//docs:spi\_vsim\_test}.
\end{itemize}

\subsection*{Used by}

Nothing in the tree yet. It is meant for the board's configuration
flash, for an SD card in SPI mode, and for sensors on the header;
issue 146 says so.

\subsection*{Limits}

One chip select, so a second chip wants a second instance or a pin
from elsewhere. One byte at a time, so a driver spends a bus
transaction per byte and the interrupt is per byte, not per command.
No quad or dual mode, and no direct memory access: reading a large
image through this costs the core a read and a write for every byte.
The board's configuration flash is reached through the
\code{STARTUPE2} primitive rather than an ordinary pin, which this
part does not do.
